\appendix{Представление графического материала}

Графический материал, выполненный на отдельных листах,
изображен на рисунках А.1--А.\arabic{числоПлакатов}.
\setcounter{числоПлакатов}{0}

\renewcommand{\thefigure}{А.\arabic{figure}} % шаблон номера для плакатов

\begin{landscape}

\begin{плакат}
    
\includegraphics[width=0.82\linewidth]{Плакат1.eps}
    \заголовок{Сведения о ВКРБ}
    \label{pl1:image}
\end{плакат}

\begin{плакат}
    
\includegraphics[width=0.82\linewidth]{Плакат2.eps}
    \заголовок{Цель и задачи разработки}
    \label{pl2:image}
\end{плакат}

\begin{плакат}
    
\includegraphics[width=0.82\linewidth]{Плакат3.eps}
    \заголовок{Файловая система FAT32}
    \label{pl3:image}
\end{плакат}

\begin{плакат}
    
\includegraphics[width=0.82\linewidth]{Плакат4.eps}
    \заголовок{Файловая система ISO9660}
    \label{pl4:image}
\end{плакат}

\begin{плакат}
	\includegraphics[width=0.82\linewidth]{Плакат5.eps}
	\заголовок{Интерфейс виртуальной файловой системы}
	\label{pl5:image}
\end{плакат}

\begin{плакат}
	\includegraphics[width=0.82\linewidth]{Плакат6.eps}
	\заголовок{Архитектура виртуальной файловой системы}
	\label{pl6:image}
\end{плакат}

\begin{плакат}
	\includegraphics[width=0.82\linewidth]{Плакат7.eps}
	\заголовок{Диаграмма классов}
	\label{pl7:image}
\end{плакат}

\begin{плакат}
	\includegraphics[width=0.82\linewidth]{Плакат8.eps}
	\заголовок{Тестирование}
	\label{pl8:image}
\end{плакат}
	
\begin{плакат}
	\includegraphics[width=0.82\linewidth]{Плакат9.eps}
	\заголовок{Заключение}
	\label{pl9:image}
\end{плакат}

\end{landscape}
